\documentclass[12pt]{article}
\usepackage{amsmath,amssymb}

\begin{document}
\section*{{2020 AMC 10A Problems/Problem 21}}
Source: AoPS Wiki

\subsection*{Solution}
First, substitute $2^{17}$ with $x$ .
Then, the given equation becomes $\frac{x^{17}+1}{x+1}=x^{16}-x^{15}+x^{14}...-x^1+x^0$ by sum of powers factorization.
Now consider only $x^{16}-x^{15}$ . This equals $x^{15}(x-1)=x^{15} \cdot (2^{17}-1)$ .
Note that $2^{17}-1$ equals $2^{16}+2^{15}+...+1$ , by difference of powers factorization (or by considering the expansion of $2^{17}=2^{16}+2^{15}+...+2^{2}+2 + 2$ ).
Thus, we can see that $x^{16}-x^{15}$ forms the sum of 17 different powers of 2.
Applying the same method to each of $x^{14}-x^{13}$ , $x^{12}-x^{11}$ , ... , $x^{2}-x^{1}$ , we can see that each of the pairs forms the sum of 17 different powers of 2. This gives us $17 \cdot 8=136$ .
But we must count also the $x^0$ term.
Thus, Our answer is $136+1=\boxed{\textbf{(C) } 137}$ . ~seanyoon777

\end{document}
